% Shared preamble of the proof overview documents  <Folder>/<Name>_overview.tex.
% Each overview starts with
%     % Shared preamble of the proof overview documents  <Folder>/<Name>_overview.tex.
% Each overview starts with
%     % Shared preamble of the proof overview documents  <Folder>/<Name>_overview.tex.
% Each overview starts with
%     \input{../../scripts/overview_preamble.tex}
% (for ThinMatching/: \input{../scripts/overview_preamble.tex}) and is compiled from its own folder by
% scripts/build_overviews.py.
\documentclass[10pt,a4paper]{article}
\usepackage[T1]{fontenc}
\usepackage[utf8]{inputenc}
\usepackage[margin=2.2cm]{geometry}
\usepackage{amsmath,amssymb,amsthm}
\usepackage{array,booktabs,longtable}
\usepackage[dvipsnames]{xcolor}
\usepackage{enumitem}
\usepackage{tikz}
\usetikzlibrary{arrows.meta,positioning,shapes.multipart,fit,calc}
\usepackage[hidelinks]{hyperref}
\setlist{itemsep=2pt,topsep=4pt}
\renewcommand{\arraystretch}{1.2}

% Lean names, files and paper references
\newcommand{\lean}[1]{\textcolor{blue!50!black}{\texttt{\detokenize{#1}}}}
\newcommand{\file}[1]{\textcolor{OliveGreen!70!black}{\texttt{\detokenize{#1}}}}
\newcommand{\paperref}[1]{\textsf{#1}}
\newcommand{\hyp}[1]{\textcolor{BrickRed}{\textbf{#1}}}

% Boxes for the dependency diagrams
\tikzset{
  res/.style={draw, rounded corners, fill=blue!6, align=center, inner sep=4pt, font=\small},
  main/.style={res, fill=blue!18, very thick},
  cited/.style={res, fill=red!10, dashed},
  shared/.style={res, fill=gray!12},
  dep/.style={-{Stealth[length=2.2mm]}, thick},
}

\newcommand{\overviewtitle}[2]{%
  \begin{center}{\Large\bfseries #1}\\[4pt]{\small Proof overview of the Lean~4/Mathlib formalisation}\\[2pt]
  {\small #2}\end{center}}

% Standard paragraph describing the paper
\newcommand{\paperversion}{\noindent\textbf{Paper.} J.~Breitung, T.~Gervens, M.~Grohe.
  \emph{Graph Similarity Meets Optimal Transport}.\\ File \file{Graph_Similarity_Meets_Optimal_Transport.pdf} in the
  repository; statement, equation, page and line numbers refer to this PDF. \file{docs/NUMBERING.md} lists the Lean
  files and declarations of all numbered statements.\par\smallskip}

% (for ThinMatching/: \input{../scripts/overview_preamble.tex}) and is compiled from its own folder by
% scripts/build_overviews.py.
\documentclass[10pt,a4paper]{article}
\usepackage[T1]{fontenc}
\usepackage[utf8]{inputenc}
\usepackage[margin=2.2cm]{geometry}
\usepackage{amsmath,amssymb,amsthm}
\usepackage{array,booktabs,longtable}
\usepackage[dvipsnames]{xcolor}
\usepackage{enumitem}
\usepackage{tikz}
\usetikzlibrary{arrows.meta,positioning,shapes.multipart,fit,calc}
\usepackage[hidelinks]{hyperref}
\setlist{itemsep=2pt,topsep=4pt}
\renewcommand{\arraystretch}{1.2}

% Lean names, files and paper references
\newcommand{\lean}[1]{\textcolor{blue!50!black}{\texttt{\detokenize{#1}}}}
\newcommand{\file}[1]{\textcolor{OliveGreen!70!black}{\texttt{\detokenize{#1}}}}
\newcommand{\paperref}[1]{\textsf{#1}}
\newcommand{\hyp}[1]{\textcolor{BrickRed}{\textbf{#1}}}

% Boxes for the dependency diagrams
\tikzset{
  res/.style={draw, rounded corners, fill=blue!6, align=center, inner sep=4pt, font=\small},
  main/.style={res, fill=blue!18, very thick},
  cited/.style={res, fill=red!10, dashed},
  shared/.style={res, fill=gray!12},
  dep/.style={-{Stealth[length=2.2mm]}, thick},
}

\newcommand{\overviewtitle}[2]{%
  \begin{center}{\Large\bfseries #1}\\[4pt]{\small Proof overview of the Lean~4/Mathlib formalisation}\\[2pt]
  {\small #2}\end{center}}

% Standard paragraph describing the paper
\newcommand{\paperversion}{\noindent\textbf{Paper.} J.~Breitung, T.~Gervens, M.~Grohe.
  \emph{Graph Similarity Meets Optimal Transport}.\\ File \file{Graph_Similarity_Meets_Optimal_Transport.pdf} in the
  repository; statement, equation, page and line numbers refer to this PDF. \file{docs/NUMBERING.md} lists the Lean
  files and declarations of all numbered statements.\par\smallskip}

% (for ThinMatching/: \input{../scripts/overview_preamble.tex}) and is compiled from its own folder by
% scripts/build_overviews.py.
\documentclass[10pt,a4paper]{article}
\usepackage[T1]{fontenc}
\usepackage[utf8]{inputenc}
\usepackage[margin=2.2cm]{geometry}
\usepackage{amsmath,amssymb,amsthm}
\usepackage{array,booktabs,longtable}
\usepackage[dvipsnames]{xcolor}
\usepackage{enumitem}
\usepackage{tikz}
\usetikzlibrary{arrows.meta,positioning,shapes.multipart,fit,calc}
\usepackage[hidelinks]{hyperref}
\setlist{itemsep=2pt,topsep=4pt}
\renewcommand{\arraystretch}{1.2}

% Lean names, files and paper references
\newcommand{\lean}[1]{\textcolor{blue!50!black}{\texttt{\detokenize{#1}}}}
\newcommand{\file}[1]{\textcolor{OliveGreen!70!black}{\texttt{\detokenize{#1}}}}
\newcommand{\paperref}[1]{\textsf{#1}}
\newcommand{\hyp}[1]{\textcolor{BrickRed}{\textbf{#1}}}

% Boxes for the dependency diagrams
\tikzset{
  res/.style={draw, rounded corners, fill=blue!6, align=center, inner sep=4pt, font=\small},
  main/.style={res, fill=blue!18, very thick},
  cited/.style={res, fill=red!10, dashed},
  shared/.style={res, fill=gray!12},
  dep/.style={-{Stealth[length=2.2mm]}, thick},
}

\newcommand{\overviewtitle}[2]{%
  \begin{center}{\Large\bfseries #1}\\[4pt]{\small Proof overview of the Lean~4/Mathlib formalisation}\\[2pt]
  {\small #2}\end{center}}

% Standard paragraph describing the paper
\newcommand{\paperversion}{\noindent\textbf{Paper.} J.~Breitung, T.~Gervens, M.~Grohe.
  \emph{Graph Similarity Meets Optimal Transport}.\\ File \file{Graph_Similarity_Meets_Optimal_Transport.pdf} in the
  repository; statement, equation, page and line numbers refer to this PDF. \file{docs/NUMBERING.md} lists the Lean
  files and declarations of all numbered statements.\par\smallskip}

\begin{document}
\overviewtitle{Theorem 9: upper bound for the local edit distance}{Folder \file{GsimOT/Theorem9_LocalEditUpperBound}}
\paperversion

\section{Statement}

\paragraph{Paper.} Theorem~9 (Section~3.2, p.~9, l.~297--299; proof in Appendix~D.2, pp.~28--30, l.~864--947):
\emph{For all graphs $G,H$ of the same order,
$\delta_1(G,H)\le15\sqrt{\delta_1^{\mathrm{OT}}(G,H)\cdot\ln\bigl(4/\delta_1^{\mathrm{OT}}(G,H)\bigr)}$, where the
right-hand side is defined as $0$ if $\delta_1^{\mathrm{OT}}(G,H)=0$.} This is the $\delta_1$ bound of
Theorem~2 (p.~3, l.~118--121). The remark after Theorem~9 (p.~9, l.~318--320) adds that it yields
$\delta_1\le C_\alpha(\delta_1^{\mathrm{OT}})^\alpha$ for every $\alpha<1/2$.

\paragraph{Lean.} For finite types $V,W$ and simple graphs $G$ on $V$, $H$ on $W$, assuming
\hyp{\lean{hCVZ : CVZRounding}} (Lemma~30, the cited rounding theorem of [CVZ11]):
\begin{itemize}
\item \lean{GsimOT.theorem9}: if $|V|=|W|$, then
  $\delta_1(G,H)\le15\sqrt{\delta_1^{\mathrm{OT}}(G,H)\ln(4/\delta_1^{\mathrm{OT}}(G,H))}$.
\item \lean{GsimOT.theorem9_of_isDS}: the same with $\delta_1^{\mathrm{OT}}$ replaced by $C_1(Q)/n$, for every
  $Q\in\mathrm{DS}(V,W)$.
\item \lean{GsimOT.theorem9_power}: for every $\alpha<1/2$ there is $C$ such that
  $\delta_1(G,H)\le C\,(\delta_1^{\mathrm{OT}}(G,H))^\alpha$ for all $G,H$ of the same order.
\end{itemize}
In Lean $\ln0=0$ and $4/0=0$, so the paper's convention for $\delta_1^{\mathrm{OT}}=0$ is built in. Apart from
the hypothesis \lean{hCVZ}, the statement is exactly the paper's.

\section{Where the proof lives}
\begin{longtable}{>{\raggedright\arraybackslash}p{5.2cm} >{\raggedright\arraybackslash}p{10.6cm}}
\toprule \textbf{File} & \textbf{Content} \\ \midrule\endhead \bottomrule\endfoot
\file{Lemma30_DependentRounding.lean} & Lemma~30 ([CVZ11, Thm.~1.1]) as the proposition \lean{CVZRounding}
  (not proved; a hypothesis), matchings \lean{IsMatching}. \\
\file{Lemma31_GoodMatching.lean} & $\Phi_v(M)$ (\lean{errAt}); Lemma~31 (\lean{lemma31}). \\
\file{PathCutting.lean} & The combinatorics of the proof of Proposition~32: cutting paths of a partial
  injection into blocks and closing them (\lean{PathCutting.exists_perm_of_partialInj}). \\
\file{Proposition32_Rounding.lean} & Proposition~32: deterministic part \lean{prop38_of_matching} and
  \lean{proposition32}. \\
\file{Theorem9.lean} & Proof of Theorem~9 (\lean{theorem9_of_isDS}, \lean{final_estimate},
  \lean{theorem9}) and the power bound (\lean{theorem9_power}). \\
\file{Common/Lemma29_Integrality.lean} & $\rho_H$ and Lemma~29 (shared with Proposition~8). \\
\file{Common/Birkhoff.lean} & Eq.~(18): a bijection $\sigma$ inside $\mathrm{supp}(Q)$
  (\lean{exists_bij_subset_supp}). \\
\end{longtable}

\section{Proof layout}
\begin{center}
\begin{tikzpicture}
  \node[main] (t9) at (0,0) {Theorem 9\\\lean{theorem9}};
  \node[main] (pow) at (5.2,0) {power bound, p.~9\\\lean{theorem9_power}};
  \node[res] (t9q) at (0,-2.1) {Theorem 9 for each $Q$\\\lean{theorem9_of_isDS}};
  \node[res] (fin) at (5.2,-2.1) {final estimate\\\lean{final_estimate}};
  \node[res] (p32) at (-4.6,-4.3) {Proposition 32\\\lean{proposition32}};
  \node[res] (det) at (0.8,-6.5) {deterministic part\\\lean{prop38_of_matching}};
  \node[res] (l31) at (-4.6,-6.5) {Lemma 31\\\lean{lemma31}};
  \node[shared] (l29) at (2.4,-4.3) {Lemma 29 (Common)\\\lean{lemma29_rho_le}, \lean{lemma29_dOne_eq_zero}};
  \node[shared] (eq18) at (-4.6,-2.1) {Eq.~(18) (Common)\\\lean{exists_bij_subset_supp}};
  \node[res] (cut) at (0.8,-8.7) {path cutting\\\lean{exists_perm_of_partialInj}};
  \node[cited] (l30) at (-4.6,-8.7) {Lemma 30 [CVZ11]\\\lean{CVZRounding} (hypothesis)};
  \draw[dep] (pow) -- (t9);
  \draw[dep] (t9) -- (t9q);
  \draw[dep] (t9q) -- (fin);
  \draw[dep] (t9q) -- (p32);
  \draw[dep] (t9q) -- (l29);
  \draw[dep] (t9q) -- (eq18);
  \draw[dep] (p32) -- (det);
  \draw[dep] (p32) -- (l31);
  \draw[dep] (det) -- (cut);
  \draw[dep] (det) -- (l29);
  \draw[dep] (l31) -- (l30);
\end{tikzpicture}
\end{center}
Arrows mean ``uses''. Blue: this folder; grey: \file{GsimOT/Common}; red, dashed: cited result, an explicit
hypothesis. \lean{theorem9} also uses Lemma~29 directly for the isomorphic case.

\section{The proof step by step}

\paragraph{Setting, Eq.~(18) (p.~28, l.~864--871).} The paper fixes an optimal $Q$ and $L=d_1^{\mathrm{OT}}$.
The Lean works with an arbitrary $Q\in\mathrm{DS}(V,W)$ and $L:=C_1(Q)$, so that all supported local costs are
$\le L$ (\lean{localCost_le_costOne}); the reference bijection $\sigma\subseteq\mathrm{supp}(Q)$ of Eq.~(18) is
\lean{exists_bij_subset_supp}.

\paragraph{Lemma 29 (p.~28, l.~872--886).} Shared with Proposition~8 and formalised in
\file{GsimOT/Common}: $\rho_H(\sigma(v),w)\le2L$ on $\mathrm{supp}(Q)$ (\lean{lemma29_rho_le}), and $L<1$
implies that $\sigma$ is an isomorphism (\lean{lemma29_iso}), hence $d_1=0$ (\lean{lemma29_dOne_eq_zero}).

\paragraph{$\Phi_v$ and Lemma 30 (pp.~28--29, l.~887--898).} $\Phi_v(M)$ is \lean{errAt}. Lemma~30 is
\lean{CVZRounding}: for doubly stochastic $Q$ and $\gamma\in(0,1/2]$ a finitely supported distribution on
matchings with marginals $(1-\gamma)Q$ and the tail bound~(19) for every linear statistic with coefficients in
$[0,1]$.

\paragraph{Lemma 31 (p.~29, l.~899--907).} \lean{lemma31} follows the proof: $\mathbb E[\Phi_v(M)]\le L$;
Eq.~(19) with $\mu:=L$, $\tau:=2+10\ln(4n)/(\gamma L)$ bounds each bad event by $1/(4n)$; the union bound gives
$\le1/4$; Markov's inequality gives $\Pr(n-|M|>2\gamma n)\le1/2$; so some matching in the support satisfies
both bounds.

\paragraph{Proposition 32 (pp.~29--30, l.~908--935).} \lean{prop38_of_matching} takes any matching
$M\subseteq\mathrm{supp}(Q)$ with $\Phi_v(M)\le X$ and shows $d_1\le X+(n-|M|)+n/\ell+2\ell L$. The partial
injection $\mu(v)=\sigma^{-1}(w)$ for $(v,w)\in M$ is cut into blocks of $\ell$ vertices and closed into cycles
(l.~912--921) by \lean{PathCutting.exists_perm_of_partialInj}: a permutation $\tau$ that agrees with $\mu$ outside
a set $B$ of $\le m+n/\ell$ closing edges, each spanning $<\ell$ steps of $\mu$; then $\pi:=\sigma\circ\tau$.
The claim $\rho_H(\sigma(v),\pi(v))\le2\ell L$ (l.~922--926) uses Lemma~29 along the block and the triangle
inequality for $\rho_H$ (\lean{rho_triangle}). The final estimate of $L_{P_\pi}(v,\pi(v))$ (l.~927--935)
inserts $A_H(\sigma(v),\pi(v'))$ and bounds the two sums by $\Phi_v(M)+m+n/\ell$ and $\rho_H(\sigma(v),\pi(v))$.
\lean{proposition32} combines this with Lemma~31.

\paragraph{Proof of Theorem 9 (p.~30, l.~936--947).} In \lean{theorem9_of_isDS}: if $L<1$, $G\cong H$
by Lemma~29 and $\delta_1=0$. Otherwise $\ell:=\lceil\sqrt{n/(2L)}\rceil$, $\gamma^*:=\sqrt{5\ln(4n)/n}$ and
$n/\ell+2\ell L\le2\sqrt{2nL}+2L$; if $\gamma^*\le1/2$, Proposition~32 gives Eq.~(20)
$d_1\le4L+2\sqrt{2nL}+4\sqrt{5n\ln(4n)}$, and if $\gamma^*>1/2$, Eq.~(20) follows from $d_1\le n$
(\lean{dOne_le_card}). Normalising uses that $x\mapsto\ln(4x)/x$ is decreasing for $x\ge1$
(\lean{log_four_mul_div_antitone}) with $n\ge1/f\ge1$. The final estimate
$\frac{b}{\sqrt{fh}}\le4\sqrt{f/h}+\frac{2\sqrt2}{\sqrt h}+4\sqrt5\le\frac{4+2\sqrt2}{\sqrt{\ln4}}+4\sqrt5<15$
is \lean{final_estimate}. Finally \lean{theorem9} passes from $C_1(Q)/n$ to the infimum
$\delta_1^{\mathrm{OT}}$ by continuity of $f\mapsto15\sqrt{f\ln(4/f)}$ at $f>0$ (\lean{exists_lt_of_dOTOne_lt});
the case $d_1^{\mathrm{OT}}<1$ is again the isomorphic case.

\paragraph{Power bound (p.~9, l.~320).} \lean{theorem9_power}: with $\varepsilon:=1-2\alpha>0$,
$\ln x\le x^\varepsilon/\varepsilon$ gives $\sqrt{f\ln(4/f)}\le\sqrt{4^\varepsilon/\varepsilon}\,f^\alpha$
(\lean{sqrt_mul_log_le_rpow}), so $C_\alpha:=15\sqrt{4^\varepsilon/\varepsilon}$ works.

\section{Assumptions and trusted base}
\begin{itemize}
\item \textbf{Lemma~30 is a hypothesis}: \lean{CVZRounding} is an explicit argument
  \lean{hCVZ} of \lean{theorem9}, \lean{theorem9_of_isDS}, \lean{theorem9_power}, \lean{proposition32},
  \lean{lemma31}; no \texttt{axiom} is introduced. It is the paper's statement specialised to doubly stochastic
  $Q$ (the only case used; a point of the perfect matching polytope), with ``random matching'' modelled as a
  finitely supported probability distribution on subsets of $V\times W$.
\item Otherwise only \lean{propext}, \lean{Classical.choice}, \lean{Quot.sound} (checked by
  \file{scripts/CheckAxioms.lean}).
\end{itemize}

\section{Notes on the formalisation}
\begin{itemize}
\item \textbf{Arbitrary instead of optimal $Q$}: Lemmas~29, 31 and Proposition~32 hold for every
  $Q\in\mathrm{DS}(V,W)$ whose supported local costs are $\le L$; Theorem~9 follows by continuity, since
  $d_1^{\mathrm{OT}}$ is an infimum in Lean.
\item The paper splits into $L=0$ and $L>0$ (then $L\ge1$ by Lemma~29); the Lean splits into $L<1$ and $L\ge1$.
\item The number of closing edges is bounded by counting the $\le n/\ell$ cut positions plus the $m$ path ends
  instead of $\sum_j\lceil p_j/\ell\rceil$; same bound. The bound $2(k-1)L$ is used as $\le2\ell L$, as in the
  paper's claim.
\item Lemma~31 treats $n=0$ separately.
\end{itemize}

\section{How to check this result}
\begin{itemize}
\item Build: \texttt{lake build GsimOT.Theorem9\_LocalEditUpperBound.Theorem9} (or \texttt{bash build.sh}).
\item Axioms: \texttt{\#print axioms GsimOT.theorem9} (and \texttt{GsimOT.theorem9\_power}); expected:
  \lean{propext}, \lean{Classical.choice}, \lean{Quot.sound}. \lean{CVZRounding} appears as a hypothesis in the
  type, not as an axiom.
\item \file{Theorem9_correspondence.pdf} in this folder lists every paper item with its Lean declaration.
\end{itemize}
\end{document}
